\section{Introduction}

TSI proposes that an internal transition model should represent declared
structural dependencies rather than treat all state coordinates as unrelated
numbers.  The historical standalone studies in Paper~3 evaluate sealed
structural OOD, open-loop rollout, and a same-task criterion association in a
different simulator.  They do not establish the exact factorization studied
here, estimate the quantities in the historical Paper~4 audit, or constitute
replications.  The new prospective extension in Paper~3, however, is this
paper's attribution cohort: the root seed and multiplicity family are shared,
and overlapping endpoints are not independent evidence.

The original question for this benchmark was whether a discovered structural
factorization predicts held-out interventions better than diagonal or
unstructured controls.  That wording is too strong.  The finite generator is
inside the TSI hypothesis class, and the fitted TSI predictor uses the same
transition equation.  Once the graph and six parameters are uniquely recovered,
exact prediction follows analytically.  The primary audit therefore concerns
exact representability and finite-family identifiability, not an empirical
causal effect of factorization.

Structural prediction requires a chain from raw observation through entity
discovery, cross-time identity, relations, structural variables, and a
dependency graph to prediction. This study addresses the later stages of that
chain. The entity set and its cross-time identity are supplied by the generator;
what we test is graph recovery and prediction once those are given. Discovering
entities and maintaining their identity from raw observation are outside the
estimand of this audit.

We retain the frozen comparison as a transparent record and add a clearly
labelled post-review diagnostic.  That diagnostic asks three narrower questions:
how cell-level averages conceal ties and structural zeroes; whether graph choice
has the same effect under factorized and untied dense heads; and where exactness
fails under genuine two-coordinate actions or a transition term outside the
fitted TSI family.  These diagnostics expose attribution boundaries rather than
rehabilitate a universal model-superiority claim.  A subsequent prospective
study then resolves the identified design defects with a fresh stochastic
family, a complete graph-by-head design, matched seven-parameter controls,
learned routing, genuine two-mechanism compositions, multiplicity-adjusted
world-level inference, and a second-language implementation.
